\section{Công thức toán}
\label{sec:cong-thuc}

\begin{itemize}
  \item Đối với công thức nằm cùng dòng với văn bản, bạn đặt nội dung giữa hai
    dấu \texttt{\$}. Ví dụ, \verb|$a^2 + b^2 = c^2$| sẽ hiển thị thành
    $a^2 + b^2 = c^2$.
  \item Đối với công thức đứng riêng một dòng và có đánh số tự động, bạn dùng
    môi trường \texttt{equation}:
\end{itemize}
\begin{ma}
\begin{equation}
  E = mc^2
  \label{eq:einstein}
\end{equation}
\end{ma}

\subsection{Ô sửa công thức}
Để mở trình soạn thảo công thức trực quan, bạn bấm biểu tượng ✎ trên một công
thức trong Mục lục hoặc nhấn \phim{Ctrl}+\phim{K} rồi \phim{E}:

\anh[0.92\linewidth]{popup-cong-thuc}{Ô sửa công thức: bản xem trước ở phía
trên, mã nguồn ở phía dưới, và nút Ω để mở bảng ký hiệu bên phải}

\begin{itemize}
  \item Thanh nút chức năng cung cấp sẵn phân số, căn thức, số mũ, chỉ số, tổng,
    tích phân, ngoặc và chữ thường. Nếu bạn bôi đen một đoạn trước khi bấm nút
    phân số, đoạn đó sẽ tự động được đưa lên tử số.
  \item Nút \textbf{Ω} mở bảng ký hiệu toán học. Bạn chỉ cần gõ tên để tìm
    nhanh, chẳng hạn gõ \texttt{lam} cho λ, \texttt{arrow} cho các mũi tên,
    hoặc \texttt{real} cho ℝ.
  \item Trong mục \nut{More}, bạn có thể chuyển đổi giữa môi trường
    \texttt{equation} và \texttt{align}, cũng như bật hoặc tắt chế độ đánh số.
\end{itemize}

\begin{meo}
Trong lúc gõ công thức, một số cụm ký tự sẽ tự động mở rộng thành lệnh
\LaTeX{}: \texttt{//} thành \verb|\frac{}{}|, \texttt{@a} thành \verb|\alpha|,
và \texttt{->} thành \verb|\to|. Nếu muốn giữ nguyên ký tự gốc vừa gõ, bạn chỉ
cần nhấn \phim{Ctrl}+\phim{Z} ngay sau đó.
\end{meo}

\begin{meo}
Tại bất kỳ vị trí nào trong văn bản, bạn có thể nhấn \phim{Ctrl}+\phim{K} rồi
\phim{M} để mở bảng ký hiệu ngay tại con trỏ. Nếu con trỏ đang nằm ngoài môi
trường toán, ký hiệu được chèn sẽ tự động bọc trong cặp dấu \texttt{\$…\$}.
\end{meo}

\begin{thuchanh}[Bài 5]
Mở tệp \texttt{bai-05.tex}, viết công thức định lý Pythagoras, sau đó thêm một
phân số bằng thanh nút chức năng, chèn ký hiệu α từ bảng ký hiệu, và cuối cùng
dùng nút ⋯ để chèn tham chiếu tới công thức vừa tạo.
\end{thuchanh}